% =====================================================================
% Section 1 --- Introduction
% ---------------------------------------------------------------------
% Job: transformation hook; the hidden closure choice; contributions
%      (rigidity map, closure taxonomy, 5x3 on a published impulse,
%      equivalences listed but not headlined); reading-guide box.
% Mindmap: N01, N08-N10.
% Status: STUB --- narrative outline v3 (2026-09-20).
% ---------------------------------------------------------------------
% SURVIVING TEXT from the old draft (paper/main.tex sec. 1): the Kuhnian
% cleavages / axiomatic-variation vocabulary --- reuse largely as-is, but
% keep the Kuhnian paragraph to a minimum (R2.3: reframe via domain of
% applicability; the cleavage framing is the reviewer's main objection).
% =====================================================================
\section{Introduction}\label{sec:intro}

% TODO (drafting notes):
% - Open concrete: the transformation programme (Sec. 2 gives the vector).
% - Name the hidden choice: the "multiplier" everyone quotes assumes a short
%   run and a financing regime; closure is *the* axiomatic variation.
% - Contributions: (i) the rigidity map; (ii) the closure taxonomy;
%   (iii) the 5x3 matrix on a published impulse; (iv) the equivalences
%   (listed, not headlined).
% - State upfront that the IO-CGE bridge is known (Robinson 2006; Rose 1995);
%   no bridge-discovery claim (R2.1/R2.3).
% - Reading-guide box: beginners Sections 1, 2, 5; experts Sections 3, 4, 6.

% TODO (reading-guide box --- framed environment):
% \begin{framed}
% \textbf{How to read this paper.} Beginners: Sections 1, 2 and 5.
% Experts: Sections 3, 4 and 6. Proofs are in Appendix A.
% \end{framed}